\section{Algorithmic Details}

This appendix summarizes the four implementations used in the report. All methods use 16 kHz, one-second mono WAV files, the label set $C=\{g,b,d,z\}$, and the same train/development/test roles.

\subsection{\texttt{strict\_course\_fft}}
\label{alg:strict-course}

\textbf{Assumption.} Each initial voiced consonant can be represented by a short-time FFT magnitude-spectrum prototype. A maximum-energy block approximates the local evidence for the initial consonant, and normalized correlation reduces level differences.

\textbf{Procedure.}
\begin{enumerate}[nosep]
    \item Frame each training waveform with 20 ms windows, 10 ms hops, and a 512-point FFT.
    \item Compute the normalized magnitude-spectrum descriptor $P_i^{(r)}$ for each candidate $M$-frame block.
    \item For each positive training sample, select the maximum-energy block $P_i^{(\star)}$.
    \item For each class $c$, average same-class $P_i^{(\star)}$ descriptors and normalize the result to obtain $T_c$.
    \item Select $M\in\{3,4,5,6\}$ on the development set by exact-match accuracy.
    \item Re-estimate templates on train+development using the selected block length.
    \item At test time, compute $s_{i,c}=\max_r\cos(P_i^{(r)},T_c)$ over all candidate blocks.
    \item Predict $\hat c_i=\arg\max_c s_{i,c}$.
\end{enumerate}

\subsection{\texttt{course\_filterbank\_corr}}
\label{alg:course-filterbank}

\textbf{Assumption.} The onset neighborhood is more reliable than whole-word averaging. Positive-negative template differences reduce shared spectral energy, and envelope autocorrelation helps separate sustained frication from short stop releases.

\textbf{Procedure.}
\begin{enumerate}[nosep]
    \item Frame the audio with 25 ms windows and 10 ms hops, then compute a 40-bin log-mel spectrogram.
    \item Estimate onset $o_i$ from the smoothed frame-energy curve.
    \item Extract 16-frame candidate windows using offsets $\Delta\in\{-2,0,2\}$.
    \item Split each window into three segments and concatenate raw segment means with baseline-subtracted segment means to form $\phi_i^{(\Delta)}$.
    \item Compute a 24-lag normalized autocorrelation vector $a_i^{(\Delta)}$ from the same window's energy envelope.
    \item For each class, average positive and negative training features to obtain $P_c^+$, $P_c^-$, $A_c^+$, and $A_c^-$.
    \item Score each offset with $0.65[\cos(\phi,P_c^+)-\cos(\phi,P_c^-)]+0.35[\cos(a,A_c^+)-\cos(a,A_c^-)]$.
    \item Use the maximum score over the three offsets as $s_{i,c}$, then predict $\arg\max_c s_{i,c}$.
\end{enumerate}

\subsection{\texttt{mlp\_ai}}
\label{alg:mlp-ai}

\textbf{Assumption.} MFCC means, standard deviations, delta-MFCC statistics, and coarse band energies summarize enough whole-word acoustic information for a two-layer nonlinear classifier.

\textbf{Structure.}
\[
62\rightarrow128\rightarrow64\rightarrow4.
\]
Hidden layers use ReLU activations, and the output layer uses sigmoid probabilities. Training uses Adam, early stopping, and $L_2$ regularization. The development set selects thresholds, while test-time four-way classification uses $\arg\max$.

\textbf{Procedure.}
\begin{enumerate}[nosep]
    \item Extract MFCC statistics, delta-MFCC statistics, and 10-dimensional band energy from each recording.
    \item Standardize features using the training-set mean and variance.
    \item Train the two-layer MLP with binary cross entropy over four labels.
    \item Estimate per-label thresholds on the development set while preserving exclusive decoding.
    \item On the test set, output four probabilities and choose the largest one.
\end{enumerate}

\subsection{\texttt{cnn\_ai}}
\label{alg:cnn-ai}

\textbf{Assumption.} A log-mel spectrogram can be treated as a time-frequency image, and convolutional layers can learn useful local acoustic patterns without a hand-specified onset window. The assumption requires enough data for the model to ignore irrelevant positions.

\textbf{Structure.}
\[
\begin{aligned}
&\mathrm{Conv}_{24}\rightarrow\mathrm{BN}\rightarrow\mathrm{ReLU}\rightarrow\mathrm{Pool} \\
&\rightarrow\mathrm{Conv}_{48}\rightarrow\mathrm{BN}\rightarrow\mathrm{ReLU}\rightarrow\mathrm{Pool} \\
&\rightarrow\mathrm{Conv}_{96}\rightarrow\mathrm{BN}\rightarrow\mathrm{ReLU}
\rightarrow\mathrm{GlobalAvgPool} \\
&\rightarrow\mathrm{FC}_{96}\rightarrow\mathrm{Dropout}\rightarrow\mathrm{FC}_{4}.
\end{aligned}
\]

\textbf{Procedure.}
\begin{enumerate}[nosep]
    \item Extract whole-word log-mel spectrograms and standardize them with the training-set global mean and variance.
    \item Feed the spectrograms as single-channel inputs to the three-stage convolutional network.
    \item Train with positive-weighted \texttt{BCEWithLogitsLoss} and AdamW.
    \item After each epoch, compute development macro-F1 and keep the best checkpoint.
    \item On the test set, apply sigmoid to the logits and choose the largest probability class.
\end{enumerate}
